% Intended LaTeX compiler: lualatex
\documentclass[a4paper,11pt]{article}

\usepackage{amsmath}
\usepackage{fontspec}
\usepackage{graphicx}
\usepackage{longtable}
\usepackage{wrapfig}
\usepackage{rotating}
\usepackage[normalem]{ulem}
\usepackage{capt-of}
\usepackage{hyperref}
\usepackage[danish]{babel}
\usepackage[a4paper,margin=1.8cm]{geometry}
\usepackage{amsmath,array,pgffor}
\pagestyle{empty}
\renewcommand{\arraystretch}{2.2}
\newcommand{\kort}[2]{\fbox{\parbox[c][1.5cm][c]{0.45\textwidth}{\centering\small\textbf{#1}\quad #2}}}
\newcommand{\kopier}[3]{\foreach \i in {1,...,#1}{\noindent\kort{#2}{#3}\hfill\kort{#2}{#3}\par\vspace{2mm}}}
\setlength{\parindent}{0em}
\parskip 1.5ex
\date{}
\title{Logaritmeregneregler: find begrundelsen (Simpel)}
\hypersetup{
 pdfauthor={},
 pdftitle={Logaritmeregneregler: find begrundelsen (Simpel)},
 pdfkeywords={},
 pdfsubject={},
 pdfcreator={},
 pdflang={Danish}}
\begin{document}

\maketitle
Beviset står i rigtig rækkefølge. Find et kort til hver linje, og skriv kortets nummer i tabellen. Nogle kort bruges flere gange, og nogle slet ikke.

\newpage
\section*{Ark A: Produktreglen}
\label{sec:orgdaea466}

\emph{Påstand:} For alle \(a,b>0\) gælder \(\log(a\cdot b)=\log(a)+\log(b)\).

\begin{center}
\begin{tabular}{|c|l|p{3cm}|}
\hline
Nr. & Formellinje & Kort nr.\\
\hline
1 & Lad \(x=\log(a)\) og \(y=\log(b)\) & \\
2 & Så er \(a=10^x\) og \(b=10^y\) & \\
3 & \(a\cdot b=10^x\cdot 10^y\) & \\
4 & \(a\cdot b=10^{x+y}\) & \\
5 & \(\log(a\cdot b)=\log(10^{x+y})\) & \\
6 & \(\log(a\cdot b)=x+y\) & \\
7 & \(\log(a\cdot b)=\log(a)+\log(b)\) & \\
\hline
\end{tabular}
\end{center}

\textbf{Opgaver}
\begin{enumerate}
\item Test reglen med \(a=100\) og \(b=1000\). Regn begge sider ud uden lommeregner.
\item Hvorfor kræver beviset, at \(a>0\) og \(b>0\)? Hvilken linje går galt, hvis \(a=-10\)?
\item Mads skriver \(\log(2+3)=\log(2)+\log(3)\). Vis med tal, at det er forkert. Forklar, hvorfor beviset ovenfor ikke virker for en sum.
\item Skriv med egne ord: Hvad er den vigtigste idé i beviset?
\end{enumerate}

\newpage
\section*{Ark B: Kvotientreglen}
\label{sec:orgf09e78d}

\emph{Påstand:} For alle \(a,b>0\) gælder \(\log\left(\dfrac ab\right)=\log(a)-\log(b)\).

\begin{center}
\begin{tabular}{|c|l|p{3cm}|}
\hline
Nr. & Formellinje & Kort nr.\\
\hline
1 & Lad \(x=\log(a)\) og \(y=\log(b)\) & \\
2 & Så er \(a=10^x\) og \(b=10^y\) & \\
3 & \(\dfrac ab=\dfrac{10^x}{10^y}\) & \\
4 & \(\dfrac ab=10^{x-y}\) & \\
5 & \(\log\left(\dfrac ab\right)=\log(10^{x-y})\) & \\
6 & \(\log\left(\dfrac ab\right)=x-y\) & \\
7 & \(\log\left(\dfrac ab\right)=\log(a)-\log(b)\) & \\
\hline
\end{tabular}
\end{center}

\textbf{Opgaver}
\begin{enumerate}
\item Sammenlign med ark A. Hvilke linjer er næsten ens, og hvilke kort er ændret?
\item Test reglen med \(a=1000\) og \(b=10\).
\item Hvad er forskellen på \(\log\left(\frac ab\right)\) og \(\frac{\log(a)}{\log(b)}\)? Vis med tal, at de ikke er ens.
\end{enumerate}

\newpage
\section*{Ark C: Potensreglen}
\label{sec:orge8d390a}

\emph{Påstand:} For alle \(a>0\) og alle tal \(r\) gælder \(\log(a^r)=r\cdot\log(a)\).

\begin{center}
\begin{tabular}{|c|l|p{3cm}|}
\hline
Nr. & Formellinje & Kort nr.\\
\hline
1 & Lad \(x=\log(a)\) & \\
2 & Så er \(a=10^x\) & \\
3 & \(a^r=(10^x)^r\) & \\
4 & \(a^r=10^{r\cdot x}\) & \\
5 & \(\log(a^r)=\log(10^{r\cdot x})\) & \\
6 & \(\log(a^r)=r\cdot x\) & \\
7 & \(\log(a^r)=r\cdot\log(a)\) & \\
\hline
\end{tabular}
\end{center}

\textbf{Opgaver}
\begin{enumerate}
\item Test reglen med \(a=10\), \(r=3\) og med \(a=100\), \(r=\frac12\).
\item Brug reglen til at vise, at \(\log(\sqrt a)=\frac12\log(a)\). Hvad er \(r\) her?
\item Hvorfor er \((\log(a))^r\) \emph{ikke} det samme som \(\log(a^r)\)? Find et taleksempel.
\end{enumerate}

\newpage
\section*{Ark D (ekstra): Følgeregler}
\label{sec:orgcdfc7c3}

Her bruges de tidligere beviste regler som værktøj.

\emph{Påstand 1:} \(\log(1)=0\)

\begin{center}
\begin{tabular}{|c|l|p{3cm}|}
\hline
Nr. & Formellinje & Kort nr.\\
\hline
1 & \(1=10^0\) & \\
2 & \(\log(1)=\log(10^0)\) & \\
3 & \(\log(1)=0\) & \\
\hline
\end{tabular}
\end{center}

\emph{Påstand 2:} For \(a>0\) gælder \(\log\left(\dfrac1a\right)=-\log(a)\).

\begin{center}
\begin{tabular}{|c|l|p{3cm}|}
\hline
Nr. & Formellinje & Kort nr.\\
\hline
1 & \(\log\left(\dfrac1a\right)=\log(1)-\log(a)\) & \\
2 & \(\log\left(\dfrac1a\right)=0-\log(a)\) & \\
3 & \(\log\left(\dfrac1a\right)=-\log(a)\) & \\
\hline
\end{tabular}
\end{center}

\textbf{Opgaver}
\begin{enumerate}
\item Hvilke tidligere beviste regler bruges?
\item Bevis, at \(\log(a\cdot b\cdot c)=\log(a)+\log(b)+\log(c)\) ved at bruge produktreglen to gange.
\item Skriv \(\ln\)-versionen af produktreglen. Hvilket grundtal og hvilke kort skal ændres?
\end{enumerate}
\end{document}
